% Pandoc-only preamble: abstracts injected before the main body.
% This file is NOT included in the PDF build (see Makefile).
% Pandoc skips \twocolumn[...] so the abstracts in main.tex are invisible
% to DOCX; this file surfaces them at the top of the document.

\noindent\textbf{АНОТАЦІЯ}

\noindent Послідовні рекомендаційні системи на основі самоуваги ігнорують
насичений контрольований сигнал: інтенсивність кожної взаємодії
(5-зірковий рейтинг, кількість годин гри, час перебування на
сторінці, вартість покупки). Ми висунули гіпотезу, що введення
інтенсивності в обчислення уваги повинно покращити точність top-$k$
пропорційно до того, наскільки інформативним є базовий сигнал. Ми
перевіряємо гіпотезу за допомогою трьох модифікацій SASRec
(IA-SASRec-Add, IA-SASRec-Mul, IA-SASRec-Val), що відповідають
трьом алгебраїчним позиціям, де інтенсивність може бути введена у
вираз уваги. Кожен варіант додає один скаляр $\lambda$, який
оптимізується під час навчання, на кожен шар уваги, спроектований
так, щоб при $\lambda{=}0$ точно відтворювався базовий SASRec:
модель є строгою надмножиною базової моделі. Ми проводимо оцінку на
семи наборах даних (MovieLens$\times$2, Amazon$\times$2,
Steam$\times$3) з п'ятьма ініціалізаціями на конфігурацію, за
протоколом значущості з двома критеріями та перевіркою стабільності
(попарний $t$-критерій на рівні ініціалізації та попарний бутстрап
на рівні користувача з $B{=}2{,}000$). Згідно з цим протоколом
жоден варіант не досягає стійкого покращення за метриками NDCG@10,
Recall@10 або MRR@10 на жодному з наборів. Отримана картина
перевертає гіпотезу: саме на наборі даних Steam (з найбільш
насиченим сигналом інтенсивності) спостерігаються найбільші
статистично значущі погіршення результатів (до $-9{,}5\%$ за
показником NDCG@10, $p<0{,}05$). Аналіз механізму на основі
навченого параметра $\lambda$ пояснює причину невдачі: $\lambda$
вільно адаптується на наборах даних, де інтенсивність не є
інформативною (до $0{,}07$ на ML-1M), але залишається на рівні
$0{,}55$--$0{,}99$ на Steam. Ми стверджуємо, що один скалярний
вентиль на кожен шар є занадто обмеженим засобом контролю для
інжекції інтенсивності, і пропонуємо три конкретні архітектурні
рішення, спрямовані на подолання цього обмеження.

\noindent\textbf{Ключові слова:} послідовні рекомендації;
самоувага; SASRec; інтеграція додаткової інформації; увага з
урахуванням інтенсивності; відтворюваність; негативний результат;
аналіз механізму.

\bigskip

\noindent\textbf{Abstract}

\noindent Sequential recommenders based on self-attention discard a
rich supervised signal: per-interaction intensity (a $5$-star
rating, hours played, dwell time, purchase value). We hypothesised
that re-introducing intensity into the attention computation should
improve top-$k$ accuracy in proportion to how informative the
underlying signal is. We test the hypothesis with three drop-in
modifications of SASRec (IA-SASRec-Add, IA-SASRec-Mul, IA-SASRec-Val),
corresponding to the three algebraic positions where intensity
can be threaded into the attention expression. Each variant adds a
single learnable scalar $\lambda$ per attention layer, designed so
that $\lambda{=}0$ recovers vanilla SASRec exactly: the model is a
strict super-set of the baseline. We evaluate across seven datasets
(MovieLens$\times$2, Amazon$\times$2, Steam$\times$3) with five seeds
per configuration, under a dual-criterion significance protocol with a
per-seed stability check (seed-level paired $t$ and per-user paired
bootstrap with $B{=}2{,}000$). Under this protocol, no variant achieves
a robust improvement on NDCG@10, Recall@10, or MRR@10 anywhere. The
pattern inverts the hypothesis: Steam (the dataset with the richest
intensity signal) hosts the largest significant degradations (up to
$-9.5\%$ NDCG@10, $p<0.05$). Mechanism analysis through the learned
$\lambda$ explains the failure: $\lambda$ adapts freely on datasets
where intensity is not informative (down to $0.07$ on ML-1M) but stays
at $0.55$--$0.99$ on Steam, where the variants are at their worst.
We argue that a single per-layer scalar gate is too limited a control
for intensity injection, and outline three concrete designs that target
this limitation.

\noindent\textbf{Keywords:} sequential recommendation; self-attention;
SASRec; side information fusion; intensity-aware attention;
reproducibility; negative result; mechanism analysis.

\bigskip\hrule\bigskip
